\section*{Capítulo \ref{ch:b3:clt} --- Funciones
características y teorema central del límite}

\begin{solution}{exo:b3:clt:1}
Uniforme en $\intcc{-1}1$: $\varphi(\xi) =
\frac12\int_{-1}^1\eu^{\iu\xi x}\dd x =
\frac{\sin\xi}{\xi}$ (igual a $1$ en $\xi = 0$). Exponencial
$\mathcal E(\lambda)$: $\varphi(\xi) =
\lambda\int_0^\infty\eu^{(\iu\xi - \lambda)x}\dd x =
\frac{\lambda}{\lambda - \iu\xi}$ (la primitiva se anula en $+\infty$,
pues $\operatorname{Re}(\iu\xi -
\lambda) < 0$). Poisson $\mathcal P(\lambda)$: por el teorema de
transferencia para leyes discretas,
\[
\varphi(\xi) = \sum_{k\geq0}\eu^{\iu\xi
k}\,\eu^{-\lambda}\frac{\lambda^k}{k!}
= \eu^{-\lambda}\exp\bigl(\lambda\eu^{\iu\xi}\bigr)
= \exp\bigl(\lambda(\eu^{\iu\xi} - 1)\bigr).
\]
Binomial $\mathcal B(n, p)$: una suma de $n$ variables de Bernoulli
\omterm{def:b3:probability:independence}{independientes}, cada una de
f.c.\ $1 - p + p\eu^{\iu\xi}$, de modo que
$\varphi(\xi) = \bigl(1 - p + p\eu^{\iu\xi}\bigr)^n$
(la~\cref{prop:b3:clt:cfbasics}(b)). Aditividad de Poisson: si $X
\sim \mathcal P(\lambda)$, $Y \sim \mathcal P(\mu)$ son
\omterm{def:b3:probability:independence}{independientes},
\[
\varphi_{X+Y}(\xi) = \eu^{\lambda(\eu^{\iu\xi}-1)}
\eu^{\mu(\eu^{\iu\xi}-1)} =
\eu^{(\lambda+\mu)(\eu^{\iu\xi}-1)},
\]
la f.c.\ de $\mathcal P(\lambda + \mu)$; la inyectividad
(el~\cref{thm:b3:clt:injectivity}) identifica la ley.
\end{solution}

\begin{solution}{exo:b3:clt:2}
(a) $\overline{\varphi_X(\xi)} = \E\eu^{-\iu\xi X} =
\varphi_{-X}(\xi)$. Así, $\varphi_X$ es real si y solo si
$\varphi_X = \varphi_{-X}$, si y solo si (inyectividad,
el~\cref{thm:b3:clt:injectivity}) $X$ y $-X$ tienen la misma ley. (b)
Escríbase $\varphi_X(\xi_0) = \eu^{\iu\theta}$. Entonces
\[
\E\bigl[1 - \cos(\xi_0X - \theta)\bigr]
= 1 - \operatorname{Re}\bigl(\eu^{-\iu\theta}
\varphi_X(\xi_0)\bigr) = 1 - 1 = 0 .
\]
El integrando es no negativo, de modo que $\cos(\xi_0X - \theta) = 1$
casi seguramente (una variable no negativa de
\omterm{def:b3:probability:space}{esperanza} nula se anula c.s.), es
decir, $\xi_0X - \theta \in 2\pi\Z$ c.s.: $X$ toma sus valores en la
progresión aritmética $\frac{\theta}{\xi_0} + \frac{2\pi}{\xi_0}\Z$ casi
seguramente. Si $X$ tiene \omterm{ex:b3:lebesgue:gamma}{densidad}, ese
conjunto numerable es de \omterm{def:b3:measure:lebesgueouter}{medida de Lebesgue} nula, de modo que soporta
probabilidad $0$ --- contradicción; por tanto,
$\abs{\varphi_X(\xi)} < 1$ para todo $\xi \neq 0$.
\end{solution}

\begin{solution}{exo:b3:clt:3}
Por \omterm{def:b3:probability:independence}{independencia} y
la~\cref{prop:b3:clt:cfbasics}:
\[
\varphi_{X+Y}(\xi) = \eu^{\iu m_1\xi - \sigma_1^2\xi^2/2}\,
\eu^{\iu m_2\xi - \sigma_2^2\xi^2/2}
= \eu^{\iu(m_1+m_2)\xi - (\sigma_1^2+\sigma_2^2)\xi^2/2},
\]
la f.c.\ de $\mathcal N(m_1 + m_2, \sigma_1^2 + \sigma_2^2)$; la
inyectividad concluye. La estabilidad por aplicaciones afines es la
regla afín ($aX + b \sim \mathcal N(am_1 + b,
a^2\sigma_1^2)$, admitiendo el caso degenerado $a = 0$), y la
estabilidad por sumas
\omterm{def:b3:probability:independence}{independientes} se sigue por
inducción del cálculo anterior. Para
\omterm{def:b3:clt:gaussianvector}{gaussianas} \emph{dependientes}, la
suma no tiene por qué ser \omterm{def:b3:clt:gaussianvector}{gaussiana}:
en el~\cref{exo:b3:clt:9}, $X$ y $Y =
\varepsilon X$ son \omterm{def:b3:clt:gaussianvector}{gaussianas}
estándar cada una, pero $X + Y$ se anula con probabilidad $\frac12$ sin
ser c.s.\ nula, de modo que no es
\omterm{def:b3:clt:gaussianvector}{gaussiana}.
\end{solution}

\begin{solution}{exo:b3:clt:4}
(a) \emph{Implicación directa.} Sean $t$ un punto de
\omterm{def:b3:topology:continuity}{continuidad} de $F_X$ y
$\delta > 0$. Tómense las rampas
\omterm{def:b3:topology:continuity}{continuas} $f^-$ ($= 1$ en
$\intoc{-\infty}{t-\delta}$, $0$ a partir de $t$, afín en medio) y $f^+$
($= 1$ en $\intoc{-\infty}t$, $0$ a partir de $t + \delta$, afín en
medio); entonces $f^- \leq
\mathbf 1_{\intoc{-\infty}t} \leq f^+$, de modo que
\[
\E f^-(X_n) \leq F_{X_n}(t) \leq \E f^+(X_n),
\]
y los términos extremos convergen a $\E f^\pm(X)$, a su vez encajonados
entre $F_X(t - \delta)$ y $F_X(t + \delta)$. Haciendo $n \to \infty$ y
después $\delta \to 0$, y usando la
\omterm{def:b3:topology:continuity}{continuidad} de $F_X$ en $t$:
$F_{X_n}(t) \to F_X(t)$.

\emph{Recíproco.} Sea $f$ \omterm{def:b3:topology:continuity}{continua}
y acotada, $M =
\sup\abs f$, $\varepsilon > 0$. Los puntos de
\omterm{def:b3:topology:continuity}{continuidad} de $F_X$ son densos
($F_X$ tiene a lo sumo una cantidad numerable de saltos), de modo que
elíjanse puntos de \omterm{def:b3:topology:continuity}{continuidad}
$a < b$ con $F_X(a) < \varepsilon$ y $1 - F_X(b) < \varepsilon$. En el
\omterm{def:b3:topology:compact}{compacto} $\intcc ab$, la función $f$
es uniformemente \omterm{def:b3:topology:continuity}{continua}: elíjanse
puntos de \omterm{def:b3:topology:continuity}{continuidad}
$a = t_0 < t_1 < \dots < t_m = b$ de $F_X$ con oscilación de $f$ a lo
sumo $\varepsilon$ en cada $\intoc{t_{j-1}}{t_j}$, y póngase $g = \sum_j
f(t_j)\,\mathbf 1_{\intoc{t_{j-1}}{t_j}}$. Entonces $\abs{f - g}
\leq \varepsilon$ en $\intoc ab$, $\abs g \leq M$, y para $T = X_n$ o
$X$:
\[
\bigl|\E f(T) - \E g(T)\bigr| \leq \varepsilon +
2M\bigl(F_T(a) + 1 - F_T(b)\bigr).
\]
Además, $\E g(X_n) = \sum_j f(t_j)\bigl(F_{X_n}(t_j) -
F_{X_n}(t_{j-1})\bigr) \to \E g(X)$ (una suma finita de términos
convergentes, siendo todos los $t_j$ puntos de
\omterm{def:b3:topology:continuity}{continuidad}), y
$F_{X_n}(a) \to F_X(a) < \varepsilon$, $1 - F_{X_n}(b) \to 1
- F_X(b) < \varepsilon$. Ensamblando:
$\limsup_n\abs{\E f(X_n) - \E f(X)} \leq 2\varepsilon +
8M\varepsilon$; hágase $\varepsilon \to 0$.

(b) La función de distribución de la constante $c$ es
$\mathbf 1_{\intco c\infty}$,
\omterm{def:b3:topology:continuity}{continua} salvo en $c$. Para
$\varepsilon > 0$, los puntos $c - \varepsilon$ y $c +
\frac\varepsilon2$ son puntos de
\omterm{def:b3:topology:continuity}{continuidad}, de modo que
\[
\P(\abs{X_n - c} > \varepsilon) \leq F_{X_n}(c -
\varepsilon) + 1 - F_{X_n}\Bigl(c + \frac\varepsilon2\Bigr)
\longrightarrow 0 + 1 - 1 = 0 .
\]
\end{solution}

\begin{solution}{exo:b3:clt:5}
Sea $z_n = p_n(\eu^{\iu\xi} - 1)$, de modo que $\varphi_{X_n}(\xi) =
(1 + z_n)^n$ (el~\cref{exo:b3:clt:1}) y $\abs{z_n} \leq 2p_n
\to 0$ (obsérvese $p_n = \frac{np_n}n \to 0$). Tanto $1 + z_n$ como
$\eu^{z_n}$ tienen módulo a lo sumo $1$: $\abs{1 + z_n} =
\abs{(1 - p_n) + p_n\eu^{\iu\xi}} \leq 1$ por la desigualdad triangular,
y $\abs{\eu^{z_n}} = \eu^{p_n(\cos\xi - 1)}
\leq 1$. La desigualdad telescópica $\abs{a^n - b^n} \leq
n\abs{a - b}$ (demostración del~\cref{thm:b3:clt:clt}) y la cota de la
serie de potencias $\abs{\eu^z - 1 - z} \leq
\abs z^2\eu^{\abs z}$ dan
\[
\bigl|(1 + z_n)^n - \eu^{nz_n}\bigr| \leq n\bigl|1 + z_n -
\eu^{z_n}\bigr| \leq n\,\abs{z_n}^2\,\eu^{\abs{z_n}} \leq
4\eu^2\,np_n^2 = 4\eu^2\,(np_n)\,p_n \longrightarrow 0 .
\]
Como $nz_n = np_n(\eu^{\iu\xi} - 1) \to
\lambda(\eu^{\iu\xi} - 1)$, concluimos $\varphi_{X_n}(\xi)
\to \exp\bigl(\lambda(\eu^{\iu\xi} - 1)\bigr)$ para todo $\xi$: la f.c.\
de $\mathcal P(\lambda)$, y Lévy (el~\cref{thm:b3:clt:levy}) da
$X_n \Rightarrow \mathcal
P(\lambda)$. Numéricamente: $\P\bigl(\mathcal B(100, 0.02) =
0\bigr) = 0.98^{100} = \eu^{100\ln 0.98} \approx
\eu^{-2.020} \approx 0.1326$, mientras que $\P\bigl(\mathcal P(2) =
0\bigr) = \eu^{-2} \approx 0.1353$: dos puntos porcentuales de
diferencia ya con este $n$ tan burdo.
\end{solution}

\begin{solution}{exo:b3:clt:6}
(a) Un lanzamiento tiene media $\frac72$ y varianza $\frac{35}{12}$, de
modo que $S$ tiene media $3500$, varianza
$\frac{35000}{12} \approx 2916.7$ y desviación típica $\approx 54.0$.
Por el TCL,
\[
\P(S > 3600) = \P\Bigl(\frac{S - 3500}{54.0} > 1.85\Bigr)
\approx 1 - \Phi(1.85) \approx 0.032 :
\]
alrededor de un $3\%$ de probabilidad. (b)
$S \sim \mathcal B(100, \frac12)$: media $50$, desviación típica $5$.
Con la corrección de \omterm{def:b3:topology:continuity}{continuidad},
\[
\P(45 \leq S \leq 55) \approx
\Phi\Bigl(\frac{55.5 - 50}{5}\Bigr) -
\Phi\Bigl(\frac{44.5 - 50}{5}\Bigr) = 2\Phi(1.1) - 1
\approx 0.729,
\]
frente al valor exacto $0.7287$; sin la corrección,
$2\Phi(1) - 1 \approx 0.683$, con casi cinco puntos de desviación. La
corrección importa porque $S$ es una variable reticular: el átomo
$\P(S = k)$ queda bien aproximado por la masa
\omterm{def:b3:clt:gaussianvector}{gaussiana} de
$\intcc{k - \frac12}{k + \frac12}$, y recortar el intervalo en los
enteros $45$ y $55$ desecha medio átomo en cada extremo.
\end{solution}

\begin{solution}{exo:b3:clt:7}
Las variables $g(U_k)$ son i.i.d.\ (imágenes
\omterm{def:b3:lebesgue:measurable}{medibles} de variables i.i.d.),
de cuadrado \omterm{def:b3:lebesgue:l1}{integrable}, de media $I$
(teorema de transferencia, el~\cref{exo:b3:product:9}) y varianza
$\sigma^2$. Si $\sigma > 0$, el~\cref{thm:b3:clt:clt} aplicado a ellas
es exactamente la convergencia enunciada
\[
\sqrt n\,\Bigl(\frac1n\sum_{k\leq n}g(U_k) - I\Bigr) =
\frac{\sum_{k\leq n}\bigl(g(U_k) - I\bigr)}{\sqrt n}
\Longrightarrow \mathcal N(0, \sigma^2)
\]
(si $\sigma = 0$, $g$ es c.s.\ constante y el miembro izquierdo se anula
idénticamente). Por tanto, $\P\bigl(\abs{\frac1n\sum g(U_k) - I} \leq
1.96\,\sigma/\sqrt n\bigr) \to 0.95$: la barra de error $\pm
1.96\,\sigma/\sqrt n$ ve la dimensión $d$ solo a través de la constante
$\sigma$, nunca a través de la velocidad en $n$. La regla del punto
medio con $n$ nodos en dimensión $d$ tiene paso $n^{-1/d}$ y error de
orden $n^{-2/d}$ para integrandos $\mathcal C^2$. El $n^{-1/2}$ de
\omterm{ex:b3:probability:sllnapps}{Monte Carlo} decae más deprisa que
$n^{-2/d}$ exactamente cuando $\frac12 > \frac2d$, es decir, $d > 4$: a
partir de la dimensión $5$, el muestreo aleatorio gana asintóticamente a
la malla --- la maldición de la dimensión perdona a los métodos
probabilísticos, y por eso \omterm{ex:b3:probability:sllnapps}{Monte
Carlo} reina en la integración en dimensión alta.
\end{solution}

\begin{solution}{exo:b3:clt:8}
\emph{Suma.} Para $\xi$ fijo:
\[
\bigl|\E\eu^{\iu\xi(X_n+Y_n)} -
\eu^{\iu\xi c}\,\E\eu^{\iu\xi X_n}\bigr|
= \bigl|\E\bigl[\eu^{\iu\xi X_n}\bigl(\eu^{\iu\xi Y_n} -
\eu^{\iu\xi c}\bigr)\bigr]\bigr|
\leq \E\bigl|\eu^{\iu\xi(Y_n - c)} - 1\bigr| .
\]
Sepárese según el suceso $\{\abs{Y_n - c} \leq \delta\}$: allí,
$\abs{\eu^{\iu\xi(Y_n-c)} - 1} \leq \abs\xi\,\delta$ (la cuerda es más
corta que el arco); el complementario aporta a lo sumo
$2\,\P(\abs{Y_n - c} > \delta) \to 0$. Por tanto, el $\limsup$ es
$\leq \abs\xi\,\delta$ para todo $\delta > 0$: la diferencia tiende a
$0$. Como $\E\eu^{\iu\xi X_n} \to
\varphi_X(\xi)$, se obtiene $\varphi_{X_n+Y_n}(\xi) \to
\eu^{\iu\xi c}\varphi_X(\xi) = \varphi_{X+c}(\xi)$, y Lévy
(el~\cref{thm:b3:clt:levy}) da $X_n + Y_n
\Rightarrow X + c$.

\emph{Producto.} Primero, $cX_n \Rightarrow cX$:
$\varphi_{cX_n}(\xi) = \varphi_{X_n}(c\xi) \to
\varphi_X(c\xi) = \varphi_{cX}(\xi)$. Después, $(Y_n - c)X_n
\to 0$ en probabilidad: las leyes de las $X_n$ son tensas (sus f.c.\
convergen a una f.c.; véase el paso de tensión
del~\cref{thm:b3:clt:levy}), de modo que, dado $\varepsilon > 0$, tómese
$M$ con $\P(\abs{X_n} > M) \leq \varepsilon$ para todo $n$; entonces
\[
\P\bigl(\abs{(Y_n - c)X_n} > \varepsilon\bigr) \leq
\P(\abs{X_n} > M) + \P\Bigl(\abs{Y_n - c} >
\frac{\varepsilon}{M}\Bigr) \leq \varepsilon + o(1) .
\]
Escribiendo $Y_nX_n = cX_n + (Y_n - c)X_n$ y aplicando la parte de la
suma (cuya demostración solo usó $Y_n' := (Y_n - c)X_n \to 0$ en
probabilidad, con constante $0$): $Y_nX_n \Rightarrow cX$.

\emph{Aplicación.} Por la ley fuerte de los grandes números
(el~\cref{thm:b3:probability:slln}), $\hat p_n \to p$ c.s., de modo que
por \omterm{def:b3:topology:continuity}{continuidad}
$\hat\sigma_n = \sqrt{\hat p_n(1 - \hat p_n)}
\to \sigma = \sqrt{p(1 - p)} > 0$ c.s.\ y, por tanto,
$\frac{\sigma}{\hat\sigma_n} \to 1$ en probabilidad. La regla del
producto de Slutsky eleva $\frac{S_n - np}{\sigma\sqrt n}
\Rightarrow \mathcal N(0,1)$ a $\frac{S_n -
np}{\hat\sigma_n\sqrt n} = \frac{\sigma}{\hat\sigma_n}\cdot
\frac{S_n - np}{\sigma\sqrt n} \Rightarrow \mathcal N(0,1)$: el
intervalo de confianza \emph{utilizable} $\hat p_n \pm
1.96\,\hat\sigma_n/\sqrt n$, construido solo con los datos, conserva su
nivel asintótico $95\%$.
\end{solution}

\begin{solution}{exo:b3:clt:9}
(a) Separando la \omterm{def:b3:probability:space}{esperanza} según los
dos valores de $\varepsilon$
(\omterm{def:b3:probability:independence}{independencia}): para $B$
boreliana, $\P(Y \in B) =
\frac12\P(X \in B) + \frac12\P(-X \in B) = \P(X \in B)$, pues
$-X \sim X$ ($\mathcal N(0,1)$ es simétrica): $Y \sim
\mathcal N(0,1)$. Y $\operatorname{Cov}(X, Y) =
\E[\varepsilon X^2] = \E[\varepsilon]\,\E[X^2] = 0 \cdot 1 =
0$. (b) $\abs Y = \abs X$, de modo que
$\P(\abs X \leq 1,\ \abs Y \geq 2)
= 0$ mientras que $\P(\abs X \leq 1)\,\P(\abs Y \geq 2) > 0$: no
\omterm{def:b3:probability:independence}{independientes}. Si $(X, Y)$
fuera un \omterm{def:b3:clt:gaussianvector}{vector gaussiano}, $X + Y =
(1 + \varepsilon)X$ sería una variable
\omterm{def:b3:clt:gaussianvector}{gaussiana} real
(la~\cref{def:b3:clt:gaussianvector} con $t = (1,1)$); pero
$\P(X + Y = 0) = \P(\varepsilon = -1) = \frac12$, mientras que una
variable \omterm{def:b3:clt:gaussianvector}{gaussiana} solo tiene un
átomo si es c.s.\ constante --- y $X + Y$ vale $2X \neq 0$ c.s.\ en
$\{\varepsilon = 1\}$. Contradicción: $(X, Y)$ no es
\omterm{def:b3:clt:gaussianvector}{gaussiano}. (c) Cada marginal es
\omterm{def:b3:clt:gaussianvector}{gaussiana} y la covarianza se anula
y, sin embargo, falla la
\omterm{def:b3:probability:independence}{independencia} --- porque el
\emph{par} no es conjuntamente
\omterm{def:b3:clt:gaussianvector}{gaussiano}.
El~\cref{thm:b3:clt:gaussianvector}(2) no puede debilitarse a
«marginales \omterm{def:b3:clt:gaussianvector}{gaussianas}».
\end{solution}

\begin{solution}{exo:b3:clt:10}
(a) El~\cref{exo:b3:fouriertransform:1} calcula
$\widehat{\eu^{-\abs\cdot}}(\xi) = \frac{2}{1 + \xi^2}$;
siendo ambos miembros \omterm{def:b3:lebesgue:l1}{integrables}, la
inversión de Fourier (el~\cref{thm:b3:fouriertransform:inversion}) le da
la vuelta:
\[
\int_\R\eu^{\iu\xi x}\,\frac{\dd x}{\pi(1 + x^2)} =
\eu^{-\abs\xi},
\]
que es exactamente $\varphi_X(\xi)$ para una variable de Cauchy $X$. (b)
Por \omterm{def:b3:probability:independence}{independencia},
$\varphi_{S_n}(\xi) =
\bigl(\eu^{-\abs\xi}\bigr)^n = \eu^{-n\abs\xi}$, de modo que
$\varphi_{S_n/n}(\xi) = \varphi_{S_n}(\xi/n) =
\eu^{-\abs\xi}$: la media empírica $\frac{S_n}n$ vuelve a ser de Cauchy
estándar para todo $n$ (inyectividad). La media nunca se concentra: sus
fluctuaciones en el instante $10^6$ son las de una sola observación. (c)
$\E\abs{X_1} = \frac2\pi\int_0^\infty\frac{x\,\dd x}{1 +
x^2} = +\infty$: la ley de Cauchy no es
\omterm{def:b3:lebesgue:l1}{integrable}, de modo que la ley fuerte de
los grandes números (el~\cref{thm:b3:probability:slln}) no se aplica, y
el TCL (que necesita varianza finita) menos aún. Aquí fallan
genuinamente sus conclusiones, no solo sus demostraciones. Comprobación
de consistencia: $\varphi(\xi) = \eu^{-\abs\xi}$ no es derivable en $0$,
como predice la~\cref{prop:b3:clt:cfbasics}(c) leída por
contrarrecíproco para una variable no
\omterm{def:b3:lebesgue:l1}{integrable}.
\end{solution}

\begin{solution}{exo:b3:clt:11}
(a) $\varphi_{aX_1 + bX_2}(\xi) =
\eu^{-a\abs\xi}\eu^{-b\abs\xi} = \eu^{-(a+b)\abs\xi} =
\varphi_{(a+b)X_1}(\xi)$
(\omterm{def:b3:probability:independence}{independencia} y
el~\cref{exo:b3:clt:10}); la inyectividad identifica las leyes.

(b) $\varphi_{aX_1+bX_2}(\xi) = \eu^{-a^2\xi^2/2}
\eu^{-b^2\xi^2/2} = \eu^{-(a^2+b^2)\xi^2/2}$: la ley de
$\sqrt{a^2+b^2}\,X_1$.

(c) Si $\varphi_X(\xi) = \eu^{-c\abs\xi^\alpha}$, entonces
$S_n = X_1 + \dots + X_n$ tiene $\varphi_{S_n} =
\eu^{-cn\abs\xi^\alpha}$, y $S_n/n^{1/\alpha}$ vuelve a tener
$\varphi(\xi) = \eu^{-c\abs\xi^\alpha}$: autorreproducción exacta bajo
el reescalado $n^{1/\alpha}$ --- $\sqrt n$ para la
\omterm{def:b3:clt:gaussianvector}{gaussiana} ($\alpha = 2$), el propio
$n$ para Cauchy ($\alpha = 1$, el~\cref{exo:b3:clt:10}(b)). Una suma de
variables i.i.d.\ solo puede converger (tras normalización afín) a una
ley estable bajo tales convoluciones; el TCL dice que la varianza finita
fuerza la cuenca \omterm{def:b3:clt:gaussianvector}{gaussiana}, y la
cuenca de Cauchy está reservada a leyes de colas tan pesadas que
$\E X^2 = \infty$ e incluso $\E\abs
X = \infty$ --- por ejemplo, sumas de variables de Cauchy. La
universalidad tiene varias islas, indexadas por el exponente de cola
$\alpha \in \intoc02$.
\end{solution}

\begin{solution}{exo:b3:clt:12}
(a) Los indicadores $\mathbf 1_{X_k \leq t}$ son i.i.d.\ de Bernoulli de
parámetro $p = F(t)$: su suma $nF_n(t)$ es binomial $\mathcal B(n, p)$;
la ley fuerte da $F_n(t)
\to p$ c.s., y el TCL (el~\cref{thm:b3:clt:clt}) aplicado a los mismos
indicadores (varianza $p(1-p)$) da el límite
\omterm{def:b3:clt:gaussianvector}{gaussiano} enunciado.

(b) $p(1 - p)$ es máxima en $p = \frac12$, es decir, donde
$F(t) = \frac12$: en la \emph{mediana}. Estimar probabilidades de cola
es asintóticamente fácil (varianza $\to 0$ cuando $p \to 0, 1$); la
región de la mediana carga con el mayor ruido estadístico --- la curva
empírica oscila más en su parte central.

(c) Para $F$ \omterm{def:b3:topology:continuity}{continua}, las
variables $U_k = F(X_k)$ son i.i.d.\ uniformes en $\intoo01$
(el~\cref{exo:b3:probability:1}), y la monotonía de $F$ da, escribiendo
$G_n$ para la función de distribución empírica de las $U_k$:
\[
\sup_{t\in\R}\,\abs{F_n(t) - F(t)}
= \sup_{u \in \operatorname{im}F}\,\abs{G_n(u) - u}
= \sup_{u\in\intcc01}\abs{G_n(u) - u} :
\]
la primera igualdad porque $\{X_k \leq t\} = \{U_k \leq
F(t)\}$ salvo sucesos nulos (monotonía; la desigualdad estricta solo
puede fallar en los tramos planos de $F$, donde ambos miembros no
cambian), y la segunda porque una $F$
\omterm{def:b3:topology:continuity}{continua}, que va de $0$ a $1$,
alcanza todo valor de $\intoo01$ (teorema del valor intermedio), y los
extremos no añaden nada ($G_n(0) - 0 = 0$ y $G_n(1) - 1 = 0$). El
miembro derecho solo involucra uniformes: una ley para todos los $F$ ---
de modo que una única tabla de valores críticos (la de la distribución
de Kolmogorov) contrasta \emph{cualquier} modelo
\omterm{def:b3:topology:continuity}{continuo} frente a los datos.
\end{solution}

\begin{solution}{pb:b3:clt:1}
\textbf{1.} La familia $(X_1, \dots, X_n, N_1, \dots, N_n)$ es
independiente: los dos bloques
son \omterm{def:b3:probability:independence}{independientes} entre sí
por construcción y cada bloque es i.i.d. $W_i$ es una
\omterm{def:b3:lebesgue:measurable}{función medible} únicamente de las
variables $(X_j)_{j<i}$ y $(N_j)_{j>i}$, todas distintas de $X_i$ y
$N_i$: por el principio de coaliciones
(el~\cref{thm:b3:probability:independence}), $W_i$ es
independiente del par
$(X_i, N_i)$. Las descomposiciones $H_i = W_i + \frac{X_i}{\sqrt n}$ y
$H_{i-1} = W_i +
\frac{N_i}{\sqrt n}$ son inmediatas de las definiciones: pasar de $H_i$
a $H_{i-1}$ intercambia el único sumando $X_i$ por $N_i$.

\textbf{2.} Taylor--Lagrange de orden $3$: existe $c$ entre $w$ y
$w + h$ con $f(w + h) = f(w) + f'(w)h +
\frac12f''(w)h^2 + \frac16f'''(c)h^3$, y $\abs{f'''(c)}
\leq M_3$ da la cota.

\textbf{3.} Restando los dos desarrollos en el punto base común
$w = W_i$:
\[
f(H_i) - f(H_{i-1}) = f'(W_i)\,\frac{X_i - N_i}{\sqrt n} +
\frac{f''(W_i)}{2}\,\frac{X_i^2 - N_i^2}{n} + R_i,
\qquad
\abs{R_i} \leq \frac{M_3}{6}\cdot
\frac{\abs{X_i}^3 + \abs{N_i}^3}{n^{3/2}} .
\]
Tómense \omterm{def:b3:probability:space}{esperanzas}. Por la pregunta
1, $f'(W_i)$ y $f''(W_i)$ son
\omterm{def:b3:probability:independence}{independientes} de
$(X_i, N_i)$, de modo que las
\omterm{def:b3:probability:space}{esperanzas} mixtas factorizan:
\begin{align*}
\E\Bigl[f'(W_i)\,\frac{X_i - N_i}{\sqrt n}\Bigr] &=
\E\bigl[f'(W_i)\bigr]\,\frac{\E X_i - \E N_i}{\sqrt n} = 0,
\\
\E\Bigl[f''(W_i)\,\frac{X_i^2 - N_i^2}{n}\Bigr] &=
\E\bigl[f''(W_i)\bigr]\,\frac{1 - 1}{n} = 0 :
\end{align*}
los dos primeros momentos de $X_i$ y $N_i$ \emph{coinciden}, y solo
sobrevive el resto:
\[
\bigl|\E f(H_i) - \E f(H_{i-1})\bigr| \leq \E\abs{R_i} \leq
\frac{M_3}{6}\cdot\frac{\beta + \gamma}{n^{3/2}} .
\]
El tercer momento \omterm{def:b3:clt:gaussianvector}{gaussiano}:
$\gamma = \E\abs{N_1}^3 =
2\int_0^\infty x^3\,\frac{\eu^{-x^2/2}}{\sqrt{2\pi}}\,\dd x
= \frac{2}{\sqrt{2\pi}}\int_0^\infty 2u\,\eu^{-u}\dd u =
\frac{4}{\sqrt{2\pi}} = \frac{2\sqrt2}{\sqrt\pi}$ (sustitución
$u = x^2/2$ y después $\Gamma(2) = 1$).

\textbf{4.} Telescopando $\E f(T_n) - \E f(G_n) =
\sum_{i=1}^n\bigl(\E f(H_i) - \E f(H_{i-1})\bigr)$ y aplicando la
pregunta 3 a cada uno de los $n$ términos:
\[
\bigl|\E f(T_n) - \E f(G_n)\bigr| \leq
n \cdot \frac{M_3(\beta + \gamma)}{6\,n^{3/2}} =
\frac{M_3\,(\beta + \gamma)}{6\,\sqrt n} .
\]

\textbf{5.} $G_n$ es $\mathcal N(0,1)$ \emph{exactamente} para todo $n$
(una suma normalizada de \omterm{def:b3:clt:gaussianvector}{gaussianas}
estándar \omterm{def:b3:probability:independence}{independientes},
el~\cref{exo:b3:clt:3}), de modo que $\E f(G_n) = \E f(N)$ y la pregunta
4 se lee $\abs{\E f(T_n) - \E f(N)} \leq
\frac{M_3(\beta+\gamma)}{6\sqrt n} \to 0$ para $f \in
\mathcal C^3_b$. \emph{Elevación.} Sea $f$
\omterm{def:b3:topology:continuity}{continua} acotada, $M = \sup\abs f$,
$\varepsilon > 0$. Elíjase $A
\geq 1$ con $\frac1{A^2} \leq \varepsilon$: Chebyshev con $\V(T_n) = 1$
da $\P(\abs{T_n} > A) \leq \varepsilon$ para todo $n$, y análogamente
$\P(\abs N > A) \leq \varepsilon$. Sea $\chi$ de clase
$\mathcal C^\infty$ con $\mathbf 1_{\intcc{-A}A} \leq \chi \leq
\mathbf 1_{\intcc{-A-1}{A+1}}$ (una meseta regular, construida
regularizando $\mathbf 1_{\intcc{-A-\frac12}{A+\frac12}}$,
el~\cref{thm:b3:lp:regularization}); $g = f\chi$ es
\omterm{def:b3:topology:continuity}{continua} de soporte
\omterm{def:b3:topology:compact}{compacto}, luego uniformemente
\omterm{def:b3:topology:continuity}{continua}, de modo que su
regularizada $g_\eta = g * \rho_\eta$ es $\mathcal
C^\infty$ con derivadas acotadas de todos los órdenes y
$\norm{g - g_\eta}_\infty \leq \varepsilon$ para $\eta$ suficientemente
pequeño. Para $T = T_n$ o $N$, puesto que $f = g$ en $\intcc{-A}A$ y
$\abs{f - g} \leq 2M$ en todas partes:
\[
\bigl|\E f(T) - \E g_\eta(T)\bigr| \leq
\E\abs{(f - g)(T)} + \norm{g - g_\eta}_\infty
\leq 2M\,\P(\abs T > A) + \varepsilon \leq
(2M + 1)\,\varepsilon .
\]
Combinando con $\E g_\eta(T_n) \to \E g_\eta(N)$ (la pregunta 4 se
aplica: $g_\eta \in \mathcal C^3_b$):
\[
\limsup_n\;\bigl|\E f(T_n) - \E f(N)\bigr| \leq
2(2M + 1)\,\varepsilon ,
\]
y $\varepsilon$ era arbitrario: $\E f(T_n) \to \E f(N)$ para toda $f$
\omterm{def:b3:topology:continuity}{continua} acotada, es decir,
$T_n \Rightarrow
\mathcal N(0,1)$.

\textbf{6.} La idéntica distribución entró solo a través de una frase:
«$X_i$ y $N_i$ tienen los mismos dos primeros momentos». Sean, pues,
$X_1, \dots, X_n$
\omterm{def:b3:probability:independence}{independientes}, centradas, de
varianzas $\sigma_i^2$ y terceros momentos finitos,
$s_n^2 = \sum_i\sigma_i^2 > 0$, y tómense $N_i \sim
\mathcal N(0, \sigma_i^2)$
\omterm{def:b3:probability:independence}{independientes} de todo lo
demás. Defínanse los híbridos con normalización $s_n$: $H_i =
\frac1{s_n}(\sum_{j\leq i}X_j + \sum_{j>i}N_j)$. En la $i$-ésima
sustitución, $\E X_i = \E N_i = 0$ y $\E X_i^2 = \E N_i^2
= \sigma_i^2$ vuelven a matar los términos de orden $f'$ y $f''$, y el
resto da (usando $\E\abs{N_i}^3 = \sigma_i^3\gamma$ por reescalado):
\[
\bigl|\E f(H_i) - \E f(H_{i-1})\bigr| \leq
\frac{M_3}{6\,s_n^3}\bigl(\E\abs{X_i}^3 +
\sigma_i^3\gamma\bigr) .
\]
Telescopando:
\[
\Bigl|\E f\Bigl(\frac{X_1 + \dots + X_n}{s_n}\Bigr) -
\E f(N)\Bigr| \leq \frac{M_3}{6\,s_n^3}\sum_{i=1}^n
\Bigl(\E\abs{X_i}^3 + \V(X_i)^{3/2}\,\gamma\Bigr) .
\]
Como $\sigma_i^3 = (\E X_i^2)^{3/2} \leq \E\abs{X_i}^3$ (la desigualdad
de las medias de potencias, es decir, Jensen para $t \mapsto
t^{3/2}$ aplicado a $X_i^2$), el miembro derecho es a lo sumo
$\frac{M_3(1 + \gamma)}{6}\cdot
\frac{\sum_i\E\abs{X_i}^3}{s_n^3}$: bajo la \emph{condición de Lyapunov}
$\frac1{s_n^3}\sum_i\E\abs{X_i}^3 \to 0$, las sumas normalizadas
convergen en ley a $\mathcal N(0,1)$ --- el TCL sin idéntica
distribución.

\textbf{7.} Sea $\rho \in \mathcal C^\infty_c(\intoo01)$ con
$\int\rho = 1$ y póngase $\psi(x) = \int_x^1\rho(s)\dd
s$: $\psi$ es $\mathcal C^\infty$, no creciente, $\psi = 1$ en $\R_-$,
$\psi = 0$ en $\intco1\infty$; sea $K =
\norm{\psi'''}_\infty$. Para $t \in \R$ y $\delta > 0$, defínanse
$\psi_\delta(x) = \psi\bigl(\frac{x - t}\delta\bigr)$ y
$\tilde\psi_\delta(x) = \psi\bigl(\frac{x - t}\delta +
1\bigr)$: son $\mathcal C^3_b$ con tercera derivada acotada por
$K/\delta^3$, y
\[
\mathbf 1_{\intoc{-\infty}{t-\delta}} \leq
\tilde\psi_\delta \leq \mathbf 1_{\intoc{-\infty}t} \leq
\psi_\delta \leq \mathbf 1_{\intoc{-\infty}{t+\delta}} .
\]
Cota superior: por la pregunta 4 aplicada a $\psi_\delta$ (con
$M_3 = K/\delta^3$),
\[
\P(T_n \leq t) \leq \E\psi_\delta(T_n) \leq
\E\psi_\delta(N) + \frac{K(\beta +
\gamma)}{6\,\delta^3\sqrt n}
\leq \Phi(t + \delta) + \frac{K(\beta +
\gamma)}{6\,\delta^3\sqrt n}
\leq \Phi(t) + \frac{\delta}{\sqrt{2\pi}} +
\frac{K(\beta + \gamma)}{6\,\delta^3\sqrt n},
\]
porque $\Phi$ es lipschitziana de constante $\frac1{\sqrt{2\pi}}$ (su
\omterm{ex:b3:lebesgue:gamma}{densidad} está acotada por
$\frac1{\sqrt{2\pi}}$). La cota inferior simétrica vía
$\tilde\psi_\delta$ da la estimación de dos términos
\[
\sup_{t\in\R}\,\bigl|\P(T_n \leq t) - \Phi(t)\bigr| \leq
\frac{K(\beta + \gamma)}{6}\cdot\frac{1}{\delta^3\sqrt n} +
\frac{\delta}{\sqrt{2\pi}}
\qquad(\delta > 0\ \text{arbitrario}).
\]
Los dos términos se equilibran cuando $\delta^{-3}n^{-1/2} \asymp
\delta$, es decir, $\delta = n^{-1/8}$: ambos valen entonces
$O(n^{-1/8})$, una velocidad uniforme explícita válida para todo $n$.
(La velocidad óptima de Berry--Esseen $C\beta/\sqrt n$ exige la
desigualdad de regularización de Fourier; la sustitución cambia agudeza
por elementalidad \omterm{def:b3:complete:complete}{completa}.)

\textbf{8.} Para $X_i = 2B_i - 1$ (signos equilibrados): centrada,
varianza $1$, y $\abs{X_i} = 1$, de modo que $\beta = 1$. La pregunta 7
acota entonces $\sup_t\abs{\P(\frac{S_n}{\sqrt n} \leq t) -
\Phi(t)}$ de manera explícita y uniforme para \emph{todo} $n$ finito ---
un enunciado global y no asintótico sobre la función de distribución. La
estimación local del~\cref{pb:b3:product:1}, pregunta 7, da en cambio la
asintótica exacta de un átomo individual, $\P(S_{2n} = 2k) \sim
\frac{\eu^{-k^2/n}}{\sqrt{\pi n}}$: resuelve probabilidades de tamaño
$n^{-1/2}$, muy por debajo de la resolución $n^{-1/8}$ de la pregunta 7,
pero es puntual, asintótica (sin error explícito a $n$ fijo) y está
atada a esta ley reticular concreta. Precisión local frente a
uniformidad global: los dos métodos son complementarios, y sumar la
estimación local en $k \in \intint{a\sqrt n}{b\sqrt n}$ recupera De
Moivre--Laplace en intervalos --- con velocidad más fina, pero solo para
esta ley.

\textbf{9.} El argumento de sustitución no usó nada de la ley de las
$X_i$ más allá de $\E X_i = 0$, $\E X_i^2 = 1$ y la finitud de
$\E\abs{X_i}^3$: si hubiéramos sustituido las
\omterm{def:b3:clt:gaussianvector}{gaussianas} $N_i$ por cualquier otra
familia i.i.d.\ con los mismos dos primeros momentos y tercer momento
finito, el mismo telescopaje acotaría
$\abs{\E f(\text{suma}_X) - \E f(\text{suma}_Y)}$ por $O(n^{-1/2})$ para
toda $f$ regular. Los estadísticos regulares de grandes sumas
\omterm{def:b3:probability:independence}{independientes} son, por tanto,
\emph{universales}: salvo un error cuantificado, dependen de la ley de
los sumandos solo a través de dos números. Este es el principio de
invariancia: demuéstrese un teorema límite para la ley más calculable
(la \omterm{def:b3:clt:gaussianvector}{gaussiana}, donde todo es exacto)
y transfiérase después a todas las leyes por sustitución. El mismo
esquema --- con sumas reemplazadas por funcionales más elaborados ---
gobierna la ley del semicírculo de Wigner para matrices aleatorias, la
universalidad de las raíces de polinomios aleatorios y buena parte de la
probabilidad moderna; el teorema central del límite es su primera y más
sencilla instancia.

\textbf{10.} $Z$ es una función boreliana solo de $(N_{i+1}, \dots,
N_n)$, mientras que $A = W_i + \theta h - Z$ y $h$ son funciones de las
variables restantes de la familia
independiente
$(X_1, \dots, X_n, N_1, \dots, N_n)$: por el principio de coaliciones
(el~\cref{thm:b3:probability:independence}), $Z$ es
independiente de $(A, h)$.
Como suma de las
\omterm{def:b3:probability:independence}{independientes}
$N_j/\sqrt n \sim \mathcal N(0, \frac1n)$, $Z \sim \mathcal
N(0, s^2)$ con $s^2 = \frac{n-i}n$ (el~\cref{exo:b3:clt:3}), de
\omterm{ex:b3:lebesgue:gamma}{densidad} acotada por
$\frac1{s\sqrt{2\pi}}$. La ley de $((A, h), Z)$ es el producto de las
dos leyes marginales, de modo que Tonelli (transferencia) congela el
primer bloque: con $G(a) =
\E\abs{g(a + Z)} = \int\abs{g(a + z)}\,\varphi_s(z)\,\dd z
\leq \frac{\norm g_{L^1}}{s\sqrt{2\pi}}$ para todo $a$,
\[
\E\bigl[\abs h^3\abs{g(A + Z)}\bigr] =
\E\bigl[\abs h^3\,G(A)\bigr] \leq
\frac{\norm g_{L^1}}{s\sqrt{2\pi}}\,\E\abs h^3
= \sqrt{\frac{n}{2\pi(n-i)}}\,\norm g_{L^1}\,\E\abs h^3 .
\]

\textbf{11.} La forma integral de la fórmula de Taylor se sigue
integrando $f(w + h) - f(w) =
h\int_0^1f'(w + \theta h)\,\dd\theta$ por partes dos veces en $\theta$.
Tomando \omterm{def:b3:probability:space}{esperanzas} en la $i$-ésima
sustitución, los órdenes $0, 1, 2$ se cancelan exactamente como en la
pregunta 3, y los dos restos (para $h = X_i/\sqrt n$ y $N_i/\sqrt n$) se
acotan, para $i \leq n - 1$, por la pregunta 10 con $g = f'''$:
\[
\bigl|\E f(H_i) - \E f(H_{i-1})\bigr| \leq
\int_0^1\frac{(1-\theta)^2}2\,\dd\theta\;
\sqrt{\frac{n}{2\pi(n-i)}}\,\norm{f'''}_{L^1}
\frac{\beta + \gamma}{n^{3/2}}
= \frac{\beta + \gamma}{6\,n^{3/2}}
\sqrt{\frac{n}{2\pi(n-i)}}\,\norm{f'''}_{L^1} .
\]
Sumando, con $\sum_{i=1}^{n-1}\sqrt{\frac n{n-i}} =
\sqrt n\sum_{m=1}^{n-1}m^{-1/2} \leq 2n$, y añadiendo la cota de la
pregunta 3 para la última sustitución ($i = n$, sin
\omterm{def:b3:clt:gaussianvector}{gaussiana} restante):
\[
\bigl|\E f(T_n) - \E f(G_n)\bigr| \leq
\frac{\beta + \gamma}{3\sqrt{2\pi}}\cdot
\frac{\norm{f'''}_{L^1}}{\sqrt n} +
\frac{M_3(\beta + \gamma)}{6\,n^{3/2}} .
\]
Rampas: $\psi_\delta'''(x) =
\delta^{-3}\psi'''\bigl(\frac{x - t}\delta\bigr)$, de modo que
$M_3 = K\delta^{-3}$ con $K = \norm{\psi'''}_\infty$ y
$\norm{\psi_\delta'''}_{L^1} = \delta^{-2}
\norm{\psi'''}_{L^1} = K_1\delta^{-2}$ (sustitución). El sándwich de la
pregunta 7 da entonces
\[
\sup_t\,\bigl|\P(T_n \leq t) - \Phi(t)\bigr| \leq
\frac{K_1(\beta + \gamma)}{3\sqrt{2\pi}}\cdot
\frac1{\delta^2\sqrt n} +
\frac{K(\beta + \gamma)}{6}\cdot
\frac1{\delta^3n^{3/2}} + \frac\delta{\sqrt{2\pi}} .
\]
En $\delta = n^{-1/6}$, el primer y el tercer término valen
$O(n^{-1/6})$ y el central $O(n^{-1})$: una velocidad uniforme
$O(n^{-1/6})$, estrictamente mejor que el $n^{-1/8}$ de la pregunta 7
--- la mitad \omterm{def:b3:clt:gaussianvector}{gaussiana} del híbrido
hizo la regularización adicional.

\textbf{12.} $\E N_1^3 = 0$ (integrando impar), y la integración por
partes da $\E N_1^4 = 3\,\E N_1^2 = 3$
($\int x^3\cdot x\varphi(x)\dd x =
3\int x^2\varphi$). Para $f$ de clase $\mathcal C^4$ con derivadas
acotadas, desarróllese cada sustitución hasta el cuarto orden: los
términos de tercer orden llevan el factor $\E X_i^3 - \E N_i^3 =
0$ (la \omterm{def:b3:probability:independence}{independencia} los
factoriza como en la pregunta 3), de modo que solo sobrevive el resto de
cuarto orden $\int_0^1\frac{(1-\theta)^3}6f^{(4)}(w + \theta
h)h^4\dd\theta$, con $\int_0^1
\frac{(1-\theta)^3}6\dd\theta = \frac1{24}$ y $\E h^4 =
\beta_4n^{-2}$ o $3n^{-2}$. La pregunta 10 (con $g =
f^{(4)}$) acota las sustituciones $i \leq n - 1$, y sumando como en la
pregunta 11:
\[
\bigl|\E f(T_n) - \E f(G_n)\bigr| \leq
\frac{\beta_4 + 3}{12\sqrt{2\pi}}\cdot
\frac{\norm{f^{(4)}}_{L^1}}{n} +
\frac{M_4(\beta_4 + 3)}{24\,n^2} .
\]
Con $\norm{\psi_\delta^{(4)}}_{L^1} = K_2\delta^{-3}$ y
$M_4 = K'\delta^{-4}$, la cota para funciones de distribución pasa a ser
$C\delta^{-3}n^{-1} + C'\delta^{-4}n^{-2} +
\frac\delta{\sqrt{2\pi}}$; en $\delta = n^{-1/4}$, los términos extremos
valen $O(n^{-1/4})$ y el central $O(n^{-1})$: velocidad $O(n^{-1/4})$.

\textbf{13.} Con $k$ momentos ajustados, el resto que sobrevive por
sustitución es de orden $\E\abs h^{k+1} \asymp
n^{-(k+1)/2}$; la cota de \omterm{def:b3:clt:gaussianvector}{gaussiana}
escondida carga $\norm{f^{(k+1)}}_{L^1}$ y la suma sobre las
sustituciones aporta el factor $2n$, dando
$\asymp\norm{f^{(k+1)}}_{L^1}\,
n^{-(k-1)/2}$ para $f$ regulares. Las rampas cuestan
$\norm{\psi_\delta^{(k+1)}}_{L^1} \asymp \delta^{-k}$, de modo que el
error para funciones de distribución es $\asymp
\delta^{-k}n^{-(k-1)/2} + \delta$, equilibrado en $\delta =
n^{-(k-1)/(2k+2)}$: velocidad $n^{-(k-1)/(2k+2)}$, que vale $n^{-1/6}$
para $k = 2$, $n^{-1/4}$ para $k = 3$, y tiende a $n^{-1/2}$ solo cuando
$k \to \infty$ --- pero $k \geq 4$ forzaría $\E X_1^4 = 3$ y más allá,
es decir, una ley que ya imita a la
\omterm{def:b3:clt:gaussianvector}{gaussiana}. La saturación es
estructural: la sustitución suma $n$ errores de sustitución \emph{en
valor absoluto}, renunciando a toda cancelación entre ellos. La
demostración de Fourier compara \omterm{def:b3:clt:cf}{funciones
características}, donde los errores aparecen con sus fases oscilantes;
la desigualdad de regularización de Esseen convierte $\abs{\varphi_{T_n}
- \varphi_N}$, integrada contra $\frac{\dd\xi}{\abs\xi}$, en una cota
para funciones de distribución con un coste solo logarítmico, y entrega
el $C\beta n^{-1/2}$ de Berry--Esseen con solo tres momentos. La
sustitución cambia optimalidad por robustez --- y, como muestra la Parte
VII, por portabilidad.

\textbf{14.} $\Sigma$ es simétrica semidefinida positiva; con
$\Sigma = PDP^{\mathsf T}$ ($P$ ortogonal, $D \geq 0$ diagonal,
el~\cref{exo:b3:submanifolds:8}), la simétrica $C =
P\sqrt DP^{\mathsf T}$ cumple $C^2 = \Sigma$. Para todo $t
\in \R^2$, $\langle t, CZ\rangle = \langle Ct, Z\rangle =
(Ct)_1Z^1 + (Ct)_2Z^2$ es una combinación lineal de
\omterm{def:b3:clt:gaussianvector}{gaussianas}
\omterm{def:b3:probability:independence}{independientes}, luego
\omterm{def:b3:clt:gaussianvector}{gaussiana} (el~\cref{exo:b3:clt:3}):
$N = CZ$ es un \omterm{def:b3:clt:gaussianvector}{vector gaussiano}; su
media es $0$ y su covarianza $\E[NN^{\mathsf T}] =
C\,\E[ZZ^{\mathsf T}]\,C^{\mathsf T} = CC^{\mathsf T} =
\Sigma$. Momentos: $\norm N^3 \leq (\abs{N_1} +
\abs{N_2})^3 \leq 4(\abs{N_1}^3 + \abs{N_2}^3)$ (convexidad de $x^3$ en
$\R_+$), y cada coordenada es una
\omterm{def:b3:clt:gaussianvector}{gaussiana} real con momentos de todos
los órdenes (el~\cref{exo:b3:product:10}): $\gamma' < \infty$. Por
último, cada $\langle t, G_n\rangle = \frac1{\sqrt
n}\sum_i\langle t, N_i\rangle$ es una suma normalizada de
$\mathcal N(0, t^{\mathsf T}\Sigma t)$ i.i.d., luego exactamente
$\mathcal N(0, t^{\mathsf T}\Sigma t)$: $G_n$ es un
\omterm{def:b3:clt:gaussianvector}{vector gaussiano} de media $0$ y
covarianza $\Sigma$, y su ley es $\mathcal N(0, \Sigma)$
(la~\cref{def:b3:clt:gaussianvector}: la ley queda determinada por estos
datos).

\textbf{15.} Sean $\phi(t) = f(w + th)$, $t \in \intcc01$: $\phi$ es
$\mathcal C^3$ con
\[
\phi'''(t) = \sum_{j,k,l\in\{1,2\}}\partial_{jkl}f(w +
th)\,h_jh_kh_l, \qquad \abs{\phi'''(t)} \leq
M_3\Bigl(\sum_j\abs{h_j}\Bigr)^3 = M_3(\abs{h_1} +
\abs{h_2})^3 .
\]
Taylor--Lagrange de orden $3$ para cierto $\phi$ entre $0$ y $1$ da la
primera desigualdad; Cauchy--Schwarz da
$\abs{h_1} + \abs{h_2} \leq \sqrt2\norm h$, de donde la constante
$\frac{2\sqrt2M_3}6 = \frac{\sqrt2M_3}3$.

\textbf{16.} Defínanse $H_i$ y $W_i$ como en la pregunta 1, ahora en
$\R^2$; el argumento de coaliciones no cambia. En la $i$-ésima
sustitución, los términos de primer orden dan
$\sum_j\E[\partial_jf(W_i)]\,(\E X_{i,j} - \E N_{i,j})/
\sqrt n = 0$ y los de segundo orden dan
$\frac1{2n}\sum_{j,k}\E[\partial_{jk}f(W_i)]\,(\Sigma_{jk}
- \Sigma_{jk}) = 0$: medias y covarianzas coinciden. La pregunta 15
acota los dos restos:
\[
\bigl|\E f(H_i) - \E f(H_{i-1})\bigr| \leq
\frac{\sqrt2M_3}3\cdot\frac{\E\norm{X_i}^3 +
\E\norm{N_i}^3}{n^{3/2}}
= \frac{\sqrt2M_3(\beta' + \gamma')}{3\,n^{3/2}},
\]
y telescopando sobre las $n$ sustituciones:
\[
\Bigl|\E f\Bigl(\frac{S_n}{\sqrt n}\Bigr) - \E
f(G_n)\Bigr| \leq \frac{\sqrt2\,M_3(\beta' +
\gamma')}{3\sqrt n},
\qquad G_n \sim \mathcal N(0, \Sigma)\ \text{exactamente} .
\]
Elevación: $\E\norm{T_n}^2 = \E\norm{X_1}^2 =
\operatorname{tr}\Sigma$ (los términos cruzados se anulan por
\omterm{def:b3:probability:independence}{independencia} y centrado), de
modo que $\P(\norm{T_n} > A) \leq
\operatorname{tr}\Sigma/A^2$, y análogamente para $N$: tensión. Dada una
$f$ \omterm{def:b3:topology:continuity}{continua} acotada y $\varepsilon
> 0$, multiplíquese por una meseta regular $\chi$ igual a $1$ en la bola
de radio $A$ y soportada en radio $A + 1$ (regularícese un indicador en
$\R^2$, el~\cref{thm:b3:lp:regularization}); $g = f\chi$ es
uniformemente \omterm{def:b3:topology:continuity}{continua} de soporte
\omterm{def:b3:topology:compact}{compacto}, de modo que su
regularización bidimensional $g_\eta$ es $\mathcal C^\infty$ con
derivadas acotadas de todos los órdenes y $\norm{g - g_\eta}_\infty
\leq \varepsilon$ para $\eta$ pequeño. La cadena de tres $\varepsilon$
de la pregunta 5 se transfiere entonces palabra por palabra: $\E f(T_n)
\to \E f(N)$ para toda $f \colon \R^2
\to \R$ \omterm{def:b3:topology:continuity}{continua} acotada. Este es
el~\cref{thm:b3:clt:multiclt} para $d = 2$, ahora demostrado --- la
sustitución esquiva el teorema de Lévy bidimensional que el capítulo
había dejado admitido.

\textbf{17.} Para $g \colon \R \to \R$
\omterm{def:b3:topology:continuity}{continua} acotada, la aplicación
$x \mapsto g(\langle t, x\rangle)$ es
\omterm{def:b3:topology:continuity}{continua} y acotada en $\R^2$, de
modo que la pregunta 16 da $\E g(\langle
t, T_n\rangle) \to \E g(\langle t, N\rangle)$: toda proyección converge
en distribución, y $\langle t,
N\rangle \sim \mathcal N(0, t^{\mathsf T}\Sigma t)$. (Esta es la
dirección fácil de Cramér--Wold: la convergencia conjunta implica la de
todas las imágenes lineales.) Aplicación: $V_i = (\xi_i, \xi_i^2 - 1)$
son vectores i.i.d.\ centrados ($\E\xi_1^2 = 1$), con entradas de
covarianza $\V(\xi_1) =
1$, $\operatorname{Cov}(\xi_1, \xi_1^2 - 1) = \E\xi_1^3$ y
$\V(\xi_1^2 - 1) = \E\xi_1^4 - 1$; el tercer momento
$\E\norm{V_1}^3 \leq 4\bigl(\E\abs{\xi_1}^3 +
\E\abs{\xi_1^2 - 1}^3\bigr)$ es finito cuando $\xi_1 \in
L^6$. La pregunta 16 produce el límite
\omterm{def:b3:clt:gaussianvector}{gaussiano} conjunto exhibido, y
el~\cref{thm:b3:clt:gaussianvector}(2): las dos coordenadas del límite
son \omterm{def:b3:probability:independence}{independientes} exactamente
cuando la covarianza $\E\xi_1^3$ se anula --- para leyes simétricas, la
media empírica y la varianza empírica se desacoplan asintóticamente.

\textbf{18.} Escríbase $g(x) - g(\theta) = (g'(\theta) +
\eta(x))(x - \theta)$, donde $\eta(x) = \frac{g(x) -
g(\theta)}{x - \theta} - g'(\theta)$ para $x \neq \theta$ y
$\eta(\theta) = 0$: la derivabilidad en $\theta$ significa precisamente
$\eta(x) \to 0$ cuando $x \to \theta$. \emph{Paso 1:}
$\hat\theta_n \to \theta$ en probabilidad: para $\varepsilon > 0$ y
cualquier $A > 0$, a partir de cierto índice
$\varepsilon\sqrt n \geq A$, de modo que $\P(\abs{\hat\theta_n -
\theta} > \varepsilon) \leq \P(\abs{\sqrt
n(\hat\theta_n - \theta)} > A) \to \P(\sigma\abs N > A)$ (las funciones
de distribución convergen en los puntos de
\omterm{def:b3:topology:continuity}{continuidad} $\pm A$), y el miembro
derecho tiende a $0$ cuando $A \to
\infty$. \emph{Paso 2:} $\eta(\hat\theta_n) \to 0$ en probabilidad: dado
$\varepsilon' > 0$, tómese $\delta$ con $\abs\eta \leq \varepsilon'$ en
$\abs{x - \theta} \leq
\delta$; entonces $\P(\abs{\eta(\hat\theta_n)} > \varepsilon')
\leq \P(\abs{\hat\theta_n - \theta} > \delta) \to 0$. \emph{Paso 3:}
\[
\sqrt n\bigl(g(\hat\theta_n) - g(\theta)\bigr) =
g'(\theta)\,\sqrt n(\hat\theta_n - \theta) +
\eta(\hat\theta_n)\cdot\sqrt n(\hat\theta_n - \theta) .
\]
Por la regla del producto de Slutsky (el~\cref{exo:b3:clt:8}, con la
sucesión $\eta(\hat\theta_n) \to 0$ en probabilidad y la convergente en
ley $\sqrt n(\hat\theta_n - \theta)$), el segundo término converge en
ley a $0\cdot\mathcal N(0,
\sigma^2) = 0$, luego a $0$ en probabilidad (el~\cref{exo:b3:clt:4}(b));
el primero converge en ley a $g'(\theta)\mathcal N(0, \sigma^2)$
(Slutsky de nuevo, o la regla afín para \omterm{def:b3:clt:cf}{funciones
características}); la regla de la suma de Slutsky los ensambla: el
límite es $\mathcal N(0,
g'(\theta)^2\sigma^2)$.

\textbf{19.} (a) El TCL da $\sqrt n(\bar X_n - \mu)
\Rightarrow \mathcal N(0, \sigma^2)$; el método delta con $g(x) = x^2$,
$g'(\mu) = 2\mu$, da $\sqrt n(\bar
X_n^2 - \mu^2) \Rightarrow \mathcal N(0, 4\mu^2\sigma^2)$ --- degenerado
(límite $0$) cuando $\mu = 0$. En ese caso la fluctuación vive una
escala más arriba: $n\bar X_n^2 =
(\sqrt n\,\bar X_n)^2$, y para $t > 0$
\[
\P\bigl(n\bar X_n^2 \leq t\bigr) = \P\bigl(-\sqrt t \leq
\sqrt n\,\bar X_n \leq \sqrt t\bigr) \longrightarrow
\Phi\Bigl(\frac{\sqrt t}\sigma\Bigr) -
\Phi\Bigl(-\frac{\sqrt t}\sigma\Bigr) = \P(\sigma^2N^2
\leq t) :
\]
$n\bar X_n^2 \Rightarrow \sigma^2N^2$, el cuadrado de una
\omterm{def:b3:clt:gaussianvector}{gaussiana} (una ley «ji cuadrado»)
--- cuando la primera derivada muere, el término de segundo orden de
Taylor dicta un límite no \omterm{def:b3:clt:gaussianvector}{gaussiano}.
(b) Aquí $\sqrt n(\hat p_n - p)
\Rightarrow \mathcal N(0, p(1 - p))$ y $g(p) =
\arcsin\sqrt p$ tiene $g'(p) = \frac1{2\sqrt{p(1 - p)}}$, de modo que
$g'(p)^2\,p(1 - p) = \frac14$: el límite es $\mathcal
N(0, \frac14)$ para todo $p \in \intoo01$. En la escala $\arcsin$, la
barra de error asintótica al $95\%$ vale $\pm
\frac{0.98}{\sqrt n}$, conocida de antemano --- mientras que en
el~\cref{ex:b3:clt:confidence} la anchura involucraba el
$\sigma = \sqrt{p(1-p)}$ desconocido, que había que acotar por el peor
caso $\frac12$ o estimar: la transformación \emph{estabiliza} la
varianza.

\textbf{20.} Sean $A^* = \{k : \mu(\{k\}) > \nu(\{k\})\}$ y
$\Delta_k = \mu(\{k\}) - \nu(\{k\})$, de modo que $\sum_k\Delta_k = 0$.
Para todo $A \subseteq \N$: $\mu(A) -
\nu(A) = \sum_{k\in A}\Delta_k \leq \sum_{k\in
A^*}\Delta_k$, con igualdad en $A = A^*$; y como las partes positiva y
negativa de $(\Delta_k)$ tienen la misma masa total,
$\sum_{A^*}\Delta_k =
\frac12\sum_k\abs{\Delta_k}$. Intercambiar $\mu, \nu$ se ocupa del
signo: $d_{\mathrm{TV}}(\mu, \nu) =
\frac12\sum_k\abs{\Delta_k}$. Acoplamiento: para todo $A$,
\[
\mu(A) - \nu(A) = \E\bigl[\mathbf 1_A(X) - \mathbf
1_A(Y)\bigr] = \E\bigl[(\mathbf 1_A(X) - \mathbf
1_A(Y))\,\mathbf 1_{X\neq Y}\bigr] \leq \P(X \neq Y),
\]
y tómese el supremo en $A$.

\textbf{21.} Las dos leyes cargan: $k = 0$: $1 - p$ frente a $\eu^{-p}$,
con $\eu^{-p} > 1 - p$; $k = 1$: $p$ frente a $p\,\eu^{-p} < p$;
$k \geq 2$: $0$ frente al resto de Poisson
$1 - \eu^{-p} - p\eu^{-p} \geq 0$. Por tanto,
\[
d_{\mathrm{TV}} = \tfrac12\bigl[(\eu^{-p} - 1 + p) + (p -
p\eu^{-p}) + (1 - \eu^{-p} - p\eu^{-p})\bigr] =
\tfrac12\bigl(2p - 2p\eu^{-p}\bigr) = p(1 - \eu^{-p}),
\]
y $1 - \eu^{-p} \leq p$ da la cota $p^2$.

\textbf{22.} Escríbanse $H_{i-1} = W_i + X_i$ y $H_i = W_i +
Y_i$ con $W_i = \sum_{j<i}Y_j + \sum_{j>i}X_j$,
independiente del par
$(X_i, Y_i)$ (coaliciones). Para $A
\subseteq \N$, condicionando a los valores, en cantidad numerable, por
\omterm{def:b3:probability:independence}{independencia},
\[
\P(H_{i-1} \in A) = \sum_{k\geq0}\P(X_i = k)\,\P(W_i + k
\in A),
\]
y análogamente para $H_i$ con $Y_i$. Restando, con $c_k
= \P(W_i + k \in A) \in \intcc01$ y $\Delta_k = \P(X_i =
k) - \P(Y_i = k)$ de suma nula:
\[
\abs{\P(H_{i-1} \in A) - \P(H_i \in A)} =
\Bigl|\sum_k\Delta_k\bigl(c_k - \tfrac12\bigr)\Bigr| \leq
\tfrac12\sum_k\abs{\Delta_k} =
d_{\mathrm{TV}}\bigl(\mathcal B(1, p_i), \mathcal
P(p_i)\bigr) .
\]
Telescopando de $H_0 = S$ a $H_n = \sum_iY_i \sim
\mathcal P(\lambda)$ (el~\cref{exo:b3:clt:1}, iterado) y usando la
pregunta 21:
\[
\abs{\P(S \in A) - \P(\mathcal P(\lambda) \in A)} \leq
\sum_{i=1}^np_i\bigl(1 - \eu^{-p_i}\bigr) \leq
\sum_{i=1}^np_i^2
\qquad\text{para todo } A :
\]
la desigualdad de Le Cam. (La cota por acoplamiento de la pregunta 20 da
una vía alternativa: acóplese cada par sobre una variable uniforme de
modo que $\P(X_i \neq Y_i) \leq p_i^2$ y acótese $\P(S \neq \sum Y_i)$;
la sustitución no necesita construcción alguna.)

\textbf{23.} (a) Con $p_i = \frac\lambda n$:
$d_{\mathrm{TV}}(\text{ley de }S, \mathcal P(\lambda))
\leq \frac{\lambda^2}n$. Esto afina el~\cref{exo:b3:clt:5} por
triplicado: un error explícito en cada $n$ finito, uniformidad sobre
todos los sucesos $A$ a la vez (no un intervalo cada vez) y ninguna
necesidad de que los $p_i$ sean iguales --- solo que $\sum_ip_i^2$ sea
pequeño, por ejemplo $\sum p_i^2 \leq
\lambda\max_ip_i$: \emph{muchos sucesos raros, ninguno dominante}. (b)
Aquí $n = 500$, $p_i = \frac1{500}$, $\lambda = 1$: el modelo de Poisson
yerra a lo sumo en $500\cdot\frac1{500^2} = 0.002$ en todo suceso; en
particular, tomando $A = \{0\}$,
\[
\P(\text{ninguna carta extraviada}) = \Bigl(1 -
\frac1{500}\Bigr)^{500},
\qquad
\Bigl|\P(\text{ninguna carta extraviada}) - \eu^{-1}\Bigr| \leq
0.002,
\]
de modo que la respuesta es $\eu^{-1} \approx 0.368$ con un $0.002$
garantizado (la discrepancia verdadera es de unos $4\cdot10^{-4}$). (c)
El problema se cierra con un método y dos regímenes. Cuando $n$
contribuciones comparables llevan cada una varianza $\frac1n$, ajustar
\emph{dos} momentos frente a la
\omterm{def:b3:clt:gaussianvector}{gaussiana} hace que los errores de
sustitución valgan $o(\frac1n)$ cada uno: las sumas se vuelven
\omterm{def:b3:clt:gaussianvector}{gaussianas} --- con Taylor como
herramienta de comparación local. Cuando $n$ contribuciones son
indicadores de probabilidad $p_i$, ajustar la \emph{media} frente a un
átomo de Poisson hace que cada sustitución cueste $p_i^2$: los recuentos
de sucesos raros se vuelven de Poisson --- con la variación total como
comparación local exacta. Los mismos híbridos, el mismo telescopio,
distinta estimación local: la sustitución es una estrategia, no un
teorema, y los límites \omterm{def:b3:clt:gaussianvector}{gaussiano} y
de Poisson son sus dos dividendos más antiguos.

\textbf{24.} El TCL da $\sqrt n(\bar X_n - \mu)
\Rightarrow \mathcal N(0, \sigma^2)$, y $g(x) = \ln x$ es derivable en
$\mu > 0$ con $g'(\mu) = \frac1\mu$: el método delta (Parte VI) produce
$\sqrt n(\ln\bar X_n - \ln\mu)
\Rightarrow \mathcal N(0, \sigma^2/\mu^2)$. Deshaciendo el intervalo
$\abs{\ln\bar X_n - \ln\mu} \leq
\frac{1.96\,\sigma}{\mu\sqrt n}$ por exponenciación:
\[
\mu \in \bar X_n\cdot
\eu^{\pm1.96\,\sigma/(\mu\sqrt n)}
\qquad\text{con probabilidad asintótica } 95\%
\]
(en la práctica, $\sigma/\mu$ se sustituye por su versión empírica,
Slutsky como en el~\cref{exo:b3:clt:8}). El intervalo multiplicativo es
el natural cuando los datos son positivos con errores proporcionales a
su tamaño --- rentas, concentraciones, semividas: magnitudes que viven
en escala logarítmica, donde los intervalos aditivos simétricos podrían
incluso cruzar el cero.

\textbf{25.} $\E[(X - p)^3] = (1-p)^3p + (-p)^3(1 - p) =
p(1-p)\bigl[(1-p)^2 - p^2\bigr] = p(1-p)(1 - 2p)$. En el análisis por
sustitución (Parte IV), el término de error dominante tras ajustar dos
momentos lleva el tercer momento \emph{con signo}: para $p < \frac12$ es
positivo (la ley se inclina a la derecha: excursiones grandes y raras
por encima de la media) y la aproximación normal desplaza
sistemáticamente la masa --- subestimando la cola izquierda, corta, y
sobrestimando la derecha --- con un error de orden $n^{-1/2}$; en $p =
\frac12$ el tercer momento se anula, la Bernoulli coincide con la
\omterm{def:b3:clt:gaussianvector}{gaussiana} hasta el tercer orden y la
velocidad mejora (la pregunta de momentos ajustados de la Parte IV).
Numéricamente: $\P(S = 0) =
0.9^{20} = 0.1216$, mientras que la
\omterm{def:b3:clt:gaussianvector}{gaussiana} $\mathcal N(2, 1.8)$ da
$\Phi\bigl(\frac{0.5 -
2}{\sqrt{1.8}}\bigr) = \Phi(-1.118) \approx 0.132$: la curva normal,
ignorante del muro en $0$ y de la asimetría hacia la derecha, pone
demasiada masa abajo --- el signo del error predicho, visible en
$n = 20$.
\end{solution}
